\documentclass[a4paper,12pt]{ctexart} % 使用 ctexart 文档类，原生支持中文

% ==================== 宏包引入 ====================
\usepackage{geometry}       % 页面边距设置
\usepackage{graphicx}       % 插入图片
\usepackage{float}          % 图片位置强制固定
\usepackage{booktabs}       % 三线表
\usepackage{amsmath}        % 数学公式
\usepackage{listings}       % 代码块
\usepackage{xcolor}         % 颜色支持
\usepackage{hyperref}       % 目录跳转和超链接
\usepackage{caption}        % 图表标题格式
\usepackage{subcaption}
\captionsetup[table]{position=bottom,justification=centering,singlelinecheck=false}


% ==================== 页面设置 ====================
\geometry{left=2.5cm, right=2.5cm, top=2.5cm, bottom=2.5cm}

% ==================== 代码高亮设置 ====================
\definecolor{codegreen}{rgb}{0,0.6,0}
\definecolor{codegray}{rgb}{0.5,0.5,0.5}
\definecolor{codepurple}{rgb}{0.58,0,0.82}
\definecolor{backcolour}{rgb}{0.95,0.95,0.92}

\lstdefinestyle{mystyle}{
    backgroundcolor=\color{backcolour},   
    commentstyle=\color{codegreen},
    keywordstyle=\color{magenta},
    numberstyle=\tiny\color{codegray},
    stringstyle=\color{codepurple},
    basicstyle=\ttfamily\footnotesize,
    breakatwhitespace=false,         
    breaklines=true,
    postbreak=\mbox{\textcolor{red}{$\hookrightarrow$}\space},
    captionpos=b,                    
    keepspaces=true,                 
    numbers=left,                    
    numbersep=5pt,                  
    showspaces=false,                
    showstringspaces=false,
    showtabs=false,                  
    tabsize=2,
    columns=flexible
}
\lstset{style=mystyle, language=Python}

% ==================== 文档信息 ====================
\title{\textbf{大模型课程第一次作业}}
\author{姓名：张宗桪 \\ 学号：2023K8009991013 \\ 专业：人工智能}
\date{\today}

% ==================== 正文开始 ====================
\begin{document}

\pagestyle{plain} % 页面顶部不显示标注，仅在底部显示页码

% 封面
\maketitle
\thispagestyle{empty} % 封面不显示页码
\newpage

% 目录
\tableofcontents
\newpage
\section{实验设计}


\subsection{数据集概况}
本次实验设置了两款数据集如下

\begin{table}[H]
    \centering
    \begin{tabular}{p{0.2\textwidth} p{0.72\textwidth}}
        \toprule
        数据集名称 & 数据内容 \\
        \midrule
        纯英文数据集 & 课程提供的英文新闻评论数据。 \\
        中英混合数据集 & 一半为新华日报官方标注数据集，另一半为上述英文新闻评论数据。 \\
        \bottomrule
    \end{tabular}
    \caption{实验所用数据集说明}
\end{table}

\subsection{核心超参数设置}

\begin{table}[H]
    \centering
    \begin{tabular}{p{0.3\textwidth} p{0.62\textwidth}}
        \toprule
        参数名称 & 参数设置 \\
        \midrule
        词表大小 & 6400 与 16000 \\
        最小词频 & 2 \\
        特殊符号 & \texttt{<|endoftext|>}、\texttt{<|im\_start|>}、\texttt{<|im\_end|>} \\

        \bottomrule
    \end{tabular}
    \caption{BPE 训练核心超参数设置}
\end{table}


\subsubsection{简要说明}
BPE 的核心思想是不断寻找语料中出现频率最高的相邻符号对，并将其合并为一个新的 Token 加入词表。该过程重复进行，直到词表大小达到预设上限，或不存在满足阈值条件的可合并符号对。通过这种方式，BPE 在控制词表规模的同时，能够较好地处理未登录词与稀有词。

本次实验中 BPE 的处理流程如下：
\begin{enumerate}
    \item \textbf{预分词}：先按空白切分，再将标点独立，避免跨词合并。
    \item \textbf{初始化词表}：将语料拆解为最基础的字符单元，构成初始词表。
    \item \textbf{统计与合并}：统计当前语料中所有相邻 Token 对的频次，选择最高频的一对，合并为新 Token，并在语料中完成对应替换。
    \item \textbf{循环迭代}：重复上述统计与合并步骤，直到词表大小达到预设 \texttt{vocab\_size}，或不存在出现频率达到 \texttt{min\_frequency} 的候选符号对。
\end{enumerate}
在超参数选择上，使用6400 和 16,000 两种词表大小设置，能够观察到不同粒度的 Token 划分对模型性能的影响。同时，设置最小词频为2，确保合并的符号对具有一定的代表性，避免过度细分导致词表过大。
\subsection{预处理与清洗}
由于本次实验使用的数据集主要是规范的新闻评论文本标注集，所以主要的清洗工作就是去掉文本中的标注，并且将文本中的英文单词和标点符号分开，避免 BPE 过程中出现跨词合并的情况。


对于中文数据，主要任务就是去掉词性标注，如处理像 [中国/ns 人民/n]nt 这种，将其转换为 [中国 人民]，对于英文数据，主要任务就是将文本中的英文单词和标点符号分开，如处理像 "This is a test." 这种，将其转换为 "This is a test ."。

\section{实验结果}


\subsection{训练结果}
选取若干代表性词语，比较其在不同数据集与词表大小设置下的编码结果。

\begin{table}[H]
    \centering
    \begin{tabular}{p{0.16\textwidth} p{0.19\textwidth} p{0.19\textwidth} p{0.19\textwidth} p{0.19\textwidth}}
        \toprule
        测试词语 & 英文-6400 & 英文-16000 & 中英混合-6400 & 中英混合-16000 \\
        \midrule
        \texttt{are} & \textbf{245} & \textbf{245} & \textbf{3917} & \textbf{3917} \\
        \texttt{China} & \textbf{456} & \textbf{456} & \textbf{4255} & \textbf{4255} \\
        \texttt{人大} & \textbf{None} & \textbf{None} & \textbf{5463} & \textbf{5463} \\
        \texttt{俄罗斯} & \textbf{None} & \textbf{None} & \textbf{5085} & \textbf{5085} \\
        \bottomrule
    \end{tabular}
    \caption{词语在不同数据集与词表大小下的编码对比}
\end{table}

\subsection{测试结果}
\subsubsection{英文}
测试输入句子为：\texttt{How are you ?}。

\begin{table}[H]
    \centering
    \begin{tabular}{p{0.2\textwidth} p{0.56\textwidth} p{0.14\textwidth}}
        \toprule
        词表配置 & Token ID 序列 & Token 数 \\
        \midrule
        英文-6400 & [2430, 245, 1014, 32] & 4 \\
        英文-16000 & [2430, 245, 1014, 32] & 4 \\
        中英混合-6400 & [39, 4082, 3917, 5614, 31] & 5 \\
        中英混合-16000 & [8410, 3917, 5614, 31] & 4 \\
        \bottomrule
    \end{tabular}
    \caption{句子在不同词表配置下的编码结果对比}
\end{table}

将对应的 Token ID 序列输入 \texttt{decode\_ids.py} 进行反向解码，结果如下。

\begin{table}[H]
    \centering
    \begin{tabular}{p{0.24\textwidth} p{0.42\textwidth} p{0.24\textwidth}}
        \toprule
        词表配置 & 输入 ID 序列 & 解码结果 \\
        \midrule
        英文-6400/16000 & [2430, 245, 1014, 32] & \texttt{How are you ?} \\
        中英混合-6400 & [39, 4082, 3917, 5614, 31] & \texttt{H ow are you ?} \\
        中英混合-16000 & [8410, 3917, 5614, 31] & \texttt{How are you ?} \\
        \bottomrule
    \end{tabular}
    \caption{不同词表配置下的解码还原结果对比}
\end{table}

\subsubsection{中文}
测试输入句子为：\texttt{张宗桪}。

\begin{table}[H]
    \centering
    \begin{tabular}{p{0.2\textwidth} p{0.56\textwidth} p{0.14\textwidth}}
        \toprule
        词表配置 & Token ID 序列 & Token 数 \\
        \midrule

        中英混合-6400 & [1188, 979, 0] & 3 \\
        中英混合-16000 & [1188, 979, 0] & 3 \\
        \bottomrule
    \end{tabular}
    \caption{句子在不同词表配置下的编码结果对比}
\end{table}

将对应的 Token ID 序列输入 \texttt{decode\_ids.py} 进行反向解码，结果如下。

\begin{table}[H]
    \centering
    \begin{tabular}{p{0.24\textwidth} p{0.42\textwidth} p{0.24\textwidth}}
        \toprule
        词表配置 & 输入 ID 序列 & 解码结果 \\
        \midrule
        中英混合 & [1188, 979, 0] & \texttt{张宗 } \\
        \bottomrule
    \end{tabular}
    \caption{不同词表配置下的解码还原结果对比}
\end{table}



\section{结果分析}
\subsection{英文测试分析}
根据表4的测试结果，针对同样的英文输入“How are you?”，在“中英混合”模型下，当词表大小为6400时，句子被切分为5个Token；而当词表大小提升至16000时，Token数下降至4个。这直观地证明了：词表大小的增加会提高文本的压缩率。较大的词表允许BPE算法记录并合并更长、更多的频繁字符对，使得原本需要多个Token拼接的词汇可以直接映射为单一Token。因此，大词表能有效缩短输入序列的长度，提高模型单次可处理的信息密度；但相应的，这也会带来Embedding层参数量的线性增长。

\subsection{中文测试分析}
BPE算法的核心优势在于其对未登录词的鲁棒性。在表6中，我使用自己的带有生僻字的姓名“张宗桪”测试时，即使某个完整的人名或生僻组合没有作为独立Token被收入词表，BPE也不会直接崩溃，而是会将其回退拆解为更细粒度的子词或单字。模型将其拆分为 [1188, 979, 0] 3个Token，这说明算法将未见过的词组退化成了单个汉字进行编码，其中ID 0 在词表中对对应<|endoftext|>。虽然这种拆分可能会导致语义信息的损失，但它保证了模型在面对新词时仍能进行处理，而不是完全无法编码。

\section{我遇到的问题}
在本次实验进行中，我主要遇到了两个方面的问题。

首先是在处理中英混合数据集时，训练脚本在特定行发生异常，切词与统计结果偏离预期。原因是源数据中潜藏了特殊的不可见排版字符（如不换行空格 U+00A0）和列表符号（如 U+00B7）。先前的基础清洗脚本仅执行了常规的去首尾空白，未能将其规范化，导致这些特殊 Unicode 字符被原样带入混合语料，干扰了下游分词器的文本匹配与词频统计。再加入新的清洗逻辑后，成功过滤掉了这些异常字符，训练过程得以顺利进行。

最后就是在全量中英混合语料上进行大词表的 BPE 训练时，系统运行内存耗尽，触发 OOM 保护机制导致进程被强制终止。原因是BPE 算法在迭代统计相邻 Token 频率时，空间复杂度极高，需要在内存中维护巨大的哈希表。叠加 Python 原生数据结构在处理海量字符串时的额外内存开销，直接突破了个人 PC 的物理硬件瓶颈。通过放弃一次性全量加载，转而随机抽取原语料库的一定比例（如 10%）构建轻量级的“实验子集”的方法规避了这个问题。
\end{document}